%=============================================================================
\section{附录：即插即用 C++ 模块}
%=============================================================================

本附录提供了笔记中所有主要算法的生产级、仅头文件（header-only）C++ 实现。每个模块均自成一体——只需将头文件放入项目中即可使用。所有模块均针对嵌入式系统设计（无动态内存分配、无异常、无 STL 容器）。

\textbf{设计理念（Design Philosophy）：}
\begin{itemize}
    \item \textbf{仅头文件（Header-only）}：无需 \texttt{.cpp} 文件。直接包含即用。
    \item \textbf{无堆分配（No heap allocation）}：所有数据均存于栈上，或通过模板进行静态尺寸确定。
    \item \textbf{极少依赖（Minimal dependencies）}：仅依赖 \texttt{<cmath>} 与 \texttt{<cstring>}，无 STL。
    \item \textbf{中断安全（ISR-safe）}：所有 \texttt{update()} 方法均可在定时器中断中安全调用。
    \item \textbf{基于模板（Template-based）}：矩阵维度为编译期常量，实现零开销抽象。
\end{itemize}

\subsection{低通滤波器（Low-Pass Filter）}

一阶 IIR 低通滤波器。一行构造，一行使用。

\begin{verbatim}
// lpf.hpp
#pragma once
#include <cmath>

class LowPassFilter {
public:
    LowPassFilter() : alpha_(1.0f), y_prev_(0.0f), initialized_(false) {}

    /// 用截止频率（Hz）和采样周期（s）进行配置
    void configure(float cutoff_hz, float dt) {
        float rc = 1.0f / (2.0f * M_PI * cutoff_hz);
        alpha_ = dt / (rc + dt);
    }

    /// 直接设置平滑因子（0 < alpha <= 1）
    void set_alpha(float alpha) { alpha_ = alpha; }

    /// 处理一个采样点，以固定采样率调用
    float update(float x) {
        if (!initialized_) {
            y_prev_ = x;
            initialized_ = true;
            return x;
        }
        y_prev_ = alpha_ * x + (1.0f - alpha_) * y_prev_;
        return y_prev_;
    }

    /// 重置滤波器状态
    void reset() { initialized_ = false; y_prev_ = 0.0f; }

    /// 获取当前滤波值，不输入新采样点
    float value() const { return y_prev_; }

private:
    float alpha_;
    float y_prev_;
    bool  initialized_;
};
\end{verbatim}

\textbf{用法（Usage）：}
\begin{verbatim}
    LowPassFilter gyro_filter;
    gyro_filter.configure(30.0f, 0.001f);  // 截止频率 30 Hz，采样频率 1 kHz

    // 在控制循环 ISR 中：
    float gyro_clean = gyro_filter.update(gyro_raw);
\end{verbatim}

\subsection{PID 控制器（PID Controller）}

功能完整的 PID，包含反计算抗积分饱和（back-calculation anti-windup）、测量值微分（derivative-on-measurement）、微分滤波和输出限幅。

\begin{verbatim}
// pid.hpp
#pragma once
#include <cmath>

struct PIDConfig {
    float kp = 0.0f;
    float ki = 0.0f;
    float kd = 0.0f;
    float dt = 0.001f;         // 采样周期（s）
    float out_min = -1e6f;     // 输出下限
    float out_max =  1e6f;     // 输出上限
    float integral_max = 1e6f; // 积分项限幅
    float d_filter_N = 10.0f;  // 微分滤波系数（越大滤波越弱）
};

class PID {
public:
    PID() = default;
    explicit PID(const PIDConfig& cfg) : cfg_(cfg) {}

    void configure(const PIDConfig& cfg) { cfg_ = cfg; }

    /// 计算 PID 输出，每个控制周期调用一次
    /// setpoint: 期望值，measurement: 实测值
    float update(float setpoint, float measurement) {
        float error = setpoint - measurement;

        // --- 比例项 ---
        float p_term = cfg_.kp * error;

        // --- 带反计算抗积分饱和的积分项 ---
        integral_ += cfg_.ki * error * cfg_.dt + anti_windup_;
        integral_ = clamp(integral_, -cfg_.integral_max, cfg_.integral_max);

        // --- 测量值微分（避免微分冲击）---
        //     带一阶低通滤波器
        float d_raw = -(measurement - meas_prev_) / cfg_.dt;
        float alpha_d = cfg_.dt * cfg_.d_filter_N
                      / (1.0f + cfg_.dt * cfg_.d_filter_N);
        d_filtered_ = alpha_d * d_raw + (1.0f - alpha_d) * d_filtered_;
        float d_term = cfg_.kd * d_filtered_;

        meas_prev_ = measurement;

        // --- 总输出 ---
        float output_raw = p_term + integral_ + d_term;
        float output = clamp(output_raw, cfg_.out_min, cfg_.out_max);

        // --- 反计算抗积分饱和 ---
        if (cfg_.ki != 0.0f) {
            anti_windup_ = (output - output_raw) * (1.0f / cfg_.ki)
                         * cfg_.dt * 0.5f;
        } else {
            anti_windup_ = 0.0f;
        }

        return output;
    }

    void reset() {
        integral_ = 0.0f;
        d_filtered_ = 0.0f;
        meas_prev_ = 0.0f;
        anti_windup_ = 0.0f;
    }

    // 访问内部状态用于调试
    float get_integral() const { return integral_; }

private:
    PIDConfig cfg_;
    float integral_    = 0.0f;
    float d_filtered_  = 0.0f;
    float meas_prev_   = 0.0f;
    float anti_windup_ = 0.0f;

    static float clamp(float v, float lo, float hi) {
        return v < lo ? lo : (v > hi ? hi : v);
    }
};
\end{verbatim}

\textbf{用法（Usage）：}
\begin{verbatim}
    PIDConfig cfg;
    cfg.kp = 6.0f;  cfg.ki = 0.3f;  cfg.kd = 3.0f;
    cfg.dt = 0.001f;
    cfg.out_min = -30000.0f;  cfg.out_max = 30000.0f;  // GM6020 输出范围
    cfg.integral_max = 10000.0f;

    PID yaw_pid(cfg);

    // 在 ISR 中：
    float torque_cmd = yaw_pid.update(target_angle, current_angle);
\end{verbatim}

\subsection{给定值斜坡（Setpoint Ramp，LPF 软启动）}

专用给定值平滑器——封装低通滤波器，实现限速的给定值过渡。

\begin{verbatim}
// ramp.hpp
#pragma once
#include "lpf.hpp"

class SetpointRamp {
public:
    /// ramp_time: 阶跃变化到达 95% 所需的近似时间（秒）
    void configure(float ramp_time, float dt) {
        float cutoff = 3.0f / (2.0f * M_PI * ramp_time);  // 3*tau ~ 95%
        filter_.configure(cutoff, dt);
    }

    float update(float raw_setpoint) {
        return filter_.update(raw_setpoint);
    }

    void reset() { filter_.reset(); }

private:
    LowPassFilter filter_;
};
\end{verbatim}

\subsection{互补滤波器（Complementary Filter，IMU 倾斜估计）}

融合加速度计（重力倾斜）与陀螺仪（角速率）。

\begin{verbatim}
// complementary_filter.hpp
#pragma once
#include <cmath>

class ComplementaryFilter {
public:
    /// alpha: 陀螺仪信任系数（典型值 0.95-0.99）
    void configure(float alpha, float dt) {
        alpha_ = alpha;
        dt_ = dt;
    }

    /// accel_angle: 由加速度计计算的倾斜角（rad）
    /// gyro_rate:   陀螺仪角速率（rad/s）
    /// 返回值:      融合后的角度估计（rad）
    float update(float accel_angle, float gyro_rate) {
        angle_ = alpha_ * (angle_ + gyro_rate * dt_)
               + (1.0f - alpha_) * accel_angle;
        return angle_;
    }

    void reset(float initial_angle = 0.0f) { angle_ = initial_angle; }
    float angle() const { return angle_; }

private:
    float alpha_ = 0.98f;
    float dt_    = 0.001f;
    float angle_ = 0.0f;
};
\end{verbatim}

\subsection{轻量矩阵库（Lightweight Matrix Library）}

编译期固定尺寸矩阵，用于嵌入式卡尔曼滤波器/LQR。无堆分配，无 STL。

\begin{verbatim}
// mat.hpp
#pragma once
#include <cstring>
#include <cmath>

template<int ROWS, int COLS>
struct Mat {
    float data[ROWS][COLS] = {};

    float& operator()(int r, int c)       { return data[r][c]; }
    float  operator()(int r, int c) const { return data[r][c]; }

    static Mat zeros() { Mat m; return m; }

    static Mat identity() {
        static_assert(ROWS == COLS, "Identity requires square matrix");
        Mat m;
        for (int i = 0; i < ROWS; ++i) m.data[i][i] = 1.0f;
        return m;
    }

    // 矩阵加法
    Mat operator+(const Mat& rhs) const {
        Mat r;
        for (int i = 0; i < ROWS; ++i)
            for (int j = 0; j < COLS; ++j)
                r.data[i][j] = data[i][j] + rhs.data[i][j];
        return r;
    }

    // 矩阵减法
    Mat operator-(const Mat& rhs) const {
        Mat r;
        for (int i = 0; i < ROWS; ++i)
            for (int j = 0; j < COLS; ++j)
                r.data[i][j] = data[i][j] - rhs.data[i][j];
        return r;
    }

    // 矩阵乘法
    template<int COLS2>
    Mat<ROWS, COLS2> operator*(const Mat<COLS, COLS2>& rhs) const {
        Mat<ROWS, COLS2> r;
        for (int i = 0; i < ROWS; ++i)
            for (int j = 0; j < COLS2; ++j)
                for (int k = 0; k < COLS; ++k)
                    r.data[i][j] += data[i][k] * rhs.data[k][j];
        return r;
    }

    // 标量乘法
    Mat operator*(float s) const {
        Mat r;
        for (int i = 0; i < ROWS; ++i)
            for (int j = 0; j < COLS; ++j)
                r.data[i][j] = data[i][j] * s;
        return r;
    }

    // 转置
    Mat<COLS, ROWS> T() const {
        Mat<COLS, ROWS> r;
        for (int i = 0; i < ROWS; ++i)
            for (int j = 0; j < COLS; ++j)
                r.data[j][i] = data[i][j];
        return r;
    }

    // 2x2 或 3x3 逆矩阵（仅限小矩阵，供卡尔曼滤波器使用）
    Mat inverse() const;
};

// 1x1 逆矩阵
template<> inline Mat<1,1> Mat<1,1>::inverse() const {
    Mat<1,1> r;
    r.data[0][0] = 1.0f / data[0][0];
    return r;
}

// 2x2 逆矩阵
template<> inline Mat<2,2> Mat<2,2>::inverse() const {
    Mat<2,2> r;
    float det = data[0][0]*data[1][1] - data[0][1]*data[1][0];
    float inv_det = 1.0f / det;
    r.data[0][0] =  data[1][1] * inv_det;
    r.data[0][1] = -data[0][1] * inv_det;
    r.data[1][0] = -data[1][0] * inv_det;
    r.data[1][1] =  data[0][0] * inv_det;
    return r;
}

// 3x3 逆矩阵
template<> inline Mat<3,3> Mat<3,3>::inverse() const {
    Mat<3,3> r;
    float det = data[0][0]*(data[1][1]*data[2][2]-data[1][2]*data[2][1])
              - data[0][1]*(data[1][0]*data[2][2]-data[1][2]*data[2][0])
              + data[0][2]*(data[1][0]*data[2][1]-data[1][1]*data[2][0]);
    float inv = 1.0f / det;
    r(0,0)=(data[1][1]*data[2][2]-data[1][2]*data[2][1])*inv;
    r(0,1)=(data[0][2]*data[2][1]-data[0][1]*data[2][2])*inv;
    r(0,2)=(data[0][1]*data[1][2]-data[0][2]*data[1][1])*inv;
    r(1,0)=(data[1][2]*data[2][0]-data[1][0]*data[2][2])*inv;
    r(1,1)=(data[0][0]*data[2][2]-data[0][2]*data[2][0])*inv;
    r(1,2)=(data[0][2]*data[1][0]-data[0][0]*data[1][2])*inv;
    r(2,0)=(data[1][0]*data[2][1]-data[1][1]*data[2][0])*inv;
    r(2,1)=(data[0][1]*data[2][0]-data[0][0]*data[2][1])*inv;
    r(2,2)=(data[0][0]*data[1][1]-data[0][1]*data[1][0])*inv;
    return r;
}
\end{verbatim}

\subsection{卡尔曼滤波器（Kalman Filter）}

通用离散时间卡尔曼滤波器。模板参数在编译期确定状态维度与测量维度。

\begin{verbatim}
// kalman.hpp
#pragma once
#include "mat.hpp"

template<int N_STATE, int N_MEAS, int N_INPUT = 1>
class KalmanFilter {
public:
    using StateVec  = Mat<N_STATE, 1>;
    using MeasVec   = Mat<N_MEAS, 1>;
    using InputVec  = Mat<N_INPUT, 1>;
    using StateMat  = Mat<N_STATE, N_STATE>;
    using InputMat  = Mat<N_STATE, N_INPUT>;
    using MeasMat   = Mat<N_MEAS, N_STATE>;
    using GainMat   = Mat<N_STATE, N_MEAS>;
    using MeasCov   = Mat<N_MEAS, N_MEAS>;

    StateMat A;   // 状态转移矩阵
    InputMat B;   // 输入矩阵
    MeasMat  C;   // 测量矩阵（部分文献记为 H）
    StateMat Q;   // 过程噪声协方差
    MeasCov  R;   // 测量噪声协方差

    StateVec x;   // 状态估计
    StateMat P;   // 估计协方差

    /// 预测步骤：向前传播状态与协方差
    void predict(const InputVec& u) {
        x = A * x + B * u;
        P = A * P * A.T() + Q;
    }

    /// 预测步骤（无输入）
    void predict() {
        x = A * x;
        P = A * P * A.T() + Q;
    }

    /// 更新步骤：融合测量值
    void update(const MeasVec& z) {
        // 新息
        MeasVec y = z - C * x;
        // 新息协方差
        MeasCov S = C * P * C.T() + R;
        // 卡尔曼增益
        GainMat K = P * C.T() * S.inverse();
        // 状态更新
        x = x + K * y;
        // 协方差更新（Joseph 形式，提高数值稳定性）
        StateMat I_KC = StateMat::identity() - K * C;
        P = I_KC * P * I_KC.T() + K * R * K.T();
    }

    /// 预测 + 更新合并
    void step(const InputVec& u, const MeasVec& z) {
        predict(u);
        update(z);
    }
};
\end{verbatim}

\textbf{用法——IMU 倾斜估计（案例研究中的 2 状态卡尔曼滤波器）：}
\begin{verbatim}
    KalmanFilter<2, 1, 1> kf;
    float dt = 0.001f;

    // 状态：[角度, 陀螺仪偏置]
    kf.A(0,0) = 1.0f;  kf.A(0,1) = -dt;
    kf.A(1,0) = 0.0f;  kf.A(1,1) = 1.0f;

    kf.B(0,0) = dt;
    kf.B(1,0) = 0.0f;

    kf.C(0,0) = 1.0f;  kf.C(0,1) = 0.0f;

    // 调整这些参数
    kf.Q(0,0) = 0.001f;  kf.Q(1,1) = 0.003f;
    kf.R(0,0) = 0.03f;

    kf.P = Mat<2,2>::identity() * 1.0f;

    // 在 ISR 中：
    Mat<1,1> u;  u(0,0) = gyro_rate;
    Mat<1,1> z;  z(0,0) = accel_angle;
    kf.step(u, z);
    float estimated_angle = kf.x(0, 0);
    float estimated_bias  = kf.x(1, 0);
\end{verbatim}

\subsection{LQR 增益（离线计算辅助工具）}

LQR 增益计算需要求解 Riccati 方程，通常在 MATLAB/Python 中\textit{离线}完成。本模块在运行时应用预计算得到的增益矩阵 $K$。

\begin{verbatim}
// lqr.hpp
#pragma once
#include "mat.hpp"

template<int N_STATE, int N_INPUT>
class LQRController {
public:
    using StateVec = Mat<N_STATE, 1>;
    using InputVec = Mat<N_INPUT, 1>;
    using GainMat  = Mat<N_INPUT, N_STATE>;

    GainMat K;  // 由 MATLAB 计算后填入：K = lqr(A, B, Q, R)

    /// 计算最优控制输入：u = -K * (x - x_ref)
    InputVec compute(const StateVec& x, const StateVec& x_ref) const {
        StateVec error = x - x_ref;
        return (K * error) * (-1.0f);
    }

    /// 计算调节控制输入（x_ref = 0）
    InputVec compute(const StateVec& x) const {
        return (K * x) * (-1.0f);
    }
};
\end{verbatim}

\textbf{用法——平衡小车：}
\begin{verbatim}
    // 增益离线计算：K = lqr(A, B, diag([100,1,10,1]), 1)
    LQRController<4, 1> lqr;
    lqr.K(0,0) = -31.62f;  // 角度增益
    lqr.K(0,1) = -5.27f;   // 角速度增益
    lqr.K(0,2) = -3.16f;   // 位置增益
    lqr.K(0,3) = -4.12f;   // 速度增益

    // 在 ISR 中（状态来自卡尔曼滤波器）：
    Mat<4,1> state;
    state(0,0) = kf.x(0,0);  // 角度
    state(1,0) = kf.x(1,0);  // 角速度
    state(2,0) = encoder_pos;
    state(3,0) = encoder_vel;

    auto u = lqr.compute(state);
    set_motor_torque(u(0,0));
\end{verbatim}

\subsection{串级 PID（Cascaded PID）}

双环串级控制器（外环位置 + 内环速度），与云台案例研究中描述的结构一致。

\begin{verbatim}
// cascaded_pid.hpp
#pragma once
#include "pid.hpp"

class CascadedPID {
public:
    PID outer;   // 位置环（较慢）
    PID inner;   // 速度环（较快）

    struct Config {
        PIDConfig outer_cfg;
        PIDConfig inner_cfg;
        float ff_gain = 0.0f;  // 速度前馈增益
    };

    void configure(const Config& cfg) {
        outer.configure(cfg.outer_cfg);
        inner.configure(cfg.inner_cfg);
        ff_gain_ = cfg.ff_gain;
    }

    /// 串级完整更新
    /// pos_setpoint: 期望位置
    /// pos_measured: 实测位置
    /// vel_measured: 实测速度
    /// vel_feedforward: 前馈参考速度（可选）
    float update(float pos_setpoint, float pos_measured,
                 float vel_measured, float vel_feedforward = 0.0f) {
        // 外环：位置 -> 速度指令
        float vel_cmd = outer.update(pos_setpoint, pos_measured)
                      + ff_gain_ * vel_feedforward;

        // 内环：速度 -> 力矩/电流指令
        float output = inner.update(vel_cmd, vel_measured);
        return output;
    }

    void reset() { outer.reset(); inner.reset(); }

private:
    float ff_gain_ = 0.0f;
};
\end{verbatim}

\textbf{用法——云台 Yaw 轴：}
\begin{verbatim}
    CascadedPID::Config cfg;
    // 外环（200 Hz）
    cfg.outer_cfg = {.kp=8.0f, .ki=0.0f, .kd=0.0f, .dt=0.005f,
                     .out_min=-300.0f, .out_max=300.0f};
    // 内环（1 kHz）
    cfg.inner_cfg = {.kp=20.0f, .ki=1.0f, .kd=0.0f, .dt=0.001f,
                     .out_min=-30000.0f, .out_max=30000.0f,
                     .integral_max=10000.0f};
    cfg.ff_gain = 1.0f;

    CascadedPID gimbal;
    gimbal.configure(cfg);

    // 内环 ISR（1 kHz）——仅更新内环
    float torque = gimbal.inner.update(vel_cmd_cached, gyro_yaw);

    // 外环 ISR（200 Hz）——同时更新双环
    float torque = gimbal.update(target_angle, encoder_yaw,
                                 gyro_yaw, ref_velocity);
\end{verbatim}

\subsection{高级 PID（积分分离 + PI-D + 2自由度）}

在基础 PID 的基础上增加了第 4.5 节中所有改进特性。

\begin{verbatim}
// pid_advanced.hpp
#pragma once
#include <cmath>

struct AdvancedPIDConfig {
    float kp = 0.0f, ki = 0.0f, kd = 0.0f;
    float dt = 0.001f;
    float out_min = -1e6f, out_max = 1e6f;
    float integral_max = 1e6f;
    float d_filter_N = 10.0f;     // 微分低通滤波系数
    float setpoint_weight_b = 1.0f; // P 给定值权重（0-1）
    float setpoint_weight_c = 0.0f; // D 给定值权重（0=PI-D，1=PID）
    float integral_sep_eps = 0.0f;  // 0 = 禁用，>0 = 误差阈值
};

class AdvancedPID {
public:
    AdvancedPID() = default;
    explicit AdvancedPID(const AdvancedPIDConfig& c) : cfg_(c) {}
    void configure(const AdvancedPIDConfig& c) { cfg_ = c; }

    float update(float setpoint, float measurement) {
        float error = setpoint - measurement;

        // --- 带给定值加权的比例项（2 自由度）---
        float p_term = cfg_.kp * (cfg_.setpoint_weight_b * setpoint
                                  - measurement);

        // --- 带积分分离的积分项 ---
        if (cfg_.integral_sep_eps <= 0.0f
            || fabsf(error) < cfg_.integral_sep_eps) {
            integral_ += cfg_.ki * error * cfg_.dt;
        }
        // 抗积分饱和：反计算
        integral_ += anti_windup_;
        integral_ = clamp(integral_, -cfg_.integral_max,
                          cfg_.integral_max);

        // --- 带给定值加权 + 滤波微分的微分项 ---
        float d_input = cfg_.setpoint_weight_c * setpoint - measurement;
        float d_raw = -(d_input - d_prev_) / cfg_.dt;
        float alpha = cfg_.dt * cfg_.d_filter_N
                    / (1.0f + cfg_.dt * cfg_.d_filter_N);
        d_filtered_ = alpha * d_raw + (1.0f - alpha) * d_filtered_;
        float d_term = cfg_.kd * d_filtered_;
        d_prev_ = d_input;

        // --- 输出 ---
        float raw = p_term + integral_ + d_term;
        float out = clamp(raw, cfg_.out_min, cfg_.out_max);

        // 反计算抗积分饱和
        anti_windup_ = (cfg_.ki != 0.0f)
            ? (out - raw) / cfg_.ki * cfg_.dt * 0.5f
            : 0.0f;
        return out;
    }

    void reset() {
        integral_ = 0; d_filtered_ = 0; d_prev_ = 0; anti_windup_ = 0;
    }

private:
    AdvancedPIDConfig cfg_;
    float integral_ = 0, d_filtered_ = 0, d_prev_ = 0, anti_windup_ = 0;
    static float clamp(float v, float lo, float hi) {
        return v < lo ? lo : (v > hi ? hi : v);
    }
};
\end{verbatim}

\textbf{用法示例：}
\begin{verbatim}
    // 标准 PID（与 pid.hpp 相同）
    AdvancedPIDConfig cfg;
    cfg.kp=6; cfg.ki=0.3; cfg.kd=3; cfg.dt=0.001f;

    // PI-D（微分作用于测量值，避免微分冲击）
    cfg.setpoint_weight_c = 0.0f;  // 微分仅作用于测量值

    // I-PD（P 和 D 均作用于测量值）
    cfg.setpoint_weight_b = 0.0f;  // 比例仅作用于测量值
    cfg.setpoint_weight_c = 0.0f;  // 微分仅作用于测量值

    // 2 自由度：软化给定值响应，保持扰动抑制能力
    cfg.setpoint_weight_b = 0.7f;  // 减小给定值超调
    cfg.setpoint_weight_c = 0.0f;  // 无微分冲击

    // 积分分离：仅在接近目标时积分
    cfg.integral_sep_eps = 5.0f;   // 误差阈值（度）
\end{verbatim}

\subsection{6自由度 IMU 扩展卡尔曼滤波器（EKF for 6-DOF IMU，姿态估计）}

即用型 EKF，从 6 轴 IMU 估计横滚角、俯仰角、偏航角及陀螺仪偏置。

\begin{verbatim}
// ekf_imu.hpp — 基于 EKF 的 6 自由度 IMU 姿态估计
// 状态：[roll, pitch, yaw, bias_gx, bias_gy, bias_gz]
// 测量：加速度计 -> 横滚角、俯仰角
// 来源：控制理论实践指南
// 许可证：MIT
#pragma once
#include "mat.hpp"
#include <cmath>

class EKF_IMU {
public:
    static constexpr int NX = 6;  // 状态数
    static constexpr int NZ = 2;  // 测量数（由加速度计得到的横滚角、俯仰角）

    using State  = Mat<NX, 1>;
    using Cov    = Mat<NX, NX>;
    using MeasV  = Mat<NZ, 1>;
    using MeasC  = Mat<NZ, NZ>;
    using KGain  = Mat<NX, NZ>;
    using Hmtx   = Mat<NZ, NX>;

    State x;   // [roll, pitch, yaw, bx, by, bz]
    Cov   P;

    Cov   Q;   // 过程噪声协方差
    MeasC R;   // 测量噪声协方差

    float dt = 0.001f;  // 采样周期

    void init() {
        x = State::zeros();
        P = Cov::identity() * 0.1f;
        // 默认值（根据你的 IMU 进行调整）
        Q = Cov::zeros();
        Q(0,0)=0.001f; Q(1,1)=0.001f; Q(2,2)=0.001f;
        Q(3,3)=0.0001f; Q(4,4)=0.0001f; Q(5,5)=0.0001f;
        R(0,0) = 0.15f; R(1,1) = 0.15f;
    }

    /// 主更新函数，以 IMU 采样频率调用
    /// gx,gy,gz: 陀螺仪读数（rad/s）
    /// ax,ay,az: 加速度计读数（m/s^2）
    void update(float gx, float gy, float gz,
                float ax, float ay, float az) {
        // === 预测步骤 ===
        float phi = x(0,0), theta = x(1,0);
        float wx = gx - x(3,0), wy = gy - x(4,0), wz = gz - x(5,0);

        float sp = sinf(phi), cp = cosf(phi);
        float tt = tanf(theta), ct = cosf(theta);

        // 欧拉运动学
        float phi_dot   = wx + sp*tt*wy + cp*tt*wz;
        float theta_dot = cp*wy - sp*wz;
        float psi_dot   = (sp/ct)*wy + (cp/ct)*wz;

        // 预测状态
        State xp = x;
        xp(0,0) += phi_dot * dt;
        xp(1,0) += theta_dot * dt;
        xp(2,0) += psi_dot * dt;

        // 雅可比矩阵 F（6x6）
        Cov F = Cov::identity();
        F(0,0) += (cp*tt*wy - sp*tt*wz)*dt;
        F(0,1)  = (sp/(ct*ct)*wy + cp/(ct*ct)*wz)*dt;
        F(0,3) = -dt;  F(0,4) = -sp*tt*dt;  F(0,5) = -cp*tt*dt;
        F(1,0)  = (-sp*wy - cp*wz)*dt;
        F(1,4) = -cp*dt;  F(1,5) = sp*dt;
        F(2,0)  = (cp/ct*wy - sp/ct*wz)*dt;
        float st = sinf(theta);
        F(2,1)  = (sp*st/(ct*ct)*wy + cp*st/(ct*ct)*wz)*dt;
        F(2,4) = -sp/ct*dt;  F(2,5) = -cp/ct*dt;

        Cov Pp = F * P * F.T() + Q;

        // === 更新步骤（加速度计 -> 横滚角、俯仰角）===
        float z_roll  = atan2f(ay, az);
        float z_pitch = atan2f(-ax, sqrtf(ay*ay + az*az));

        MeasV z;  z(0,0) = z_roll;  z(1,0) = z_pitch;
        MeasV h;  h(0,0) = xp(0,0); h(1,0) = xp(1,0);

        MeasV innov;
        innov(0,0) = wrap_pi(z(0,0) - h(0,0));
        innov(1,0) = wrap_pi(z(1,0) - h(1,0));

        // H = [1 0 0 0 0 0; 0 1 0 0 0 0]
        Hmtx H = Hmtx::zeros();
        H(0,0) = 1.0f; H(1,1) = 1.0f;

        MeasC S = H * Pp * H.T() + R;
        KGain Kk = Pp * H.T() * S.inverse();

        x = xp + Kk * innov;
        P = (Cov::identity() - Kk * H) * Pp;
    }

    float roll()  const { return x(0,0); }
    float pitch() const { return x(1,0); }
    float yaw()   const { return x(2,0); }
    float bias_x() const { return x(3,0); }
    float bias_y() const { return x(4,0); }
    float bias_z() const { return x(5,0); }

private:
    static float wrap_pi(float a) {
        while (a >  3.14159265f) a -= 6.28318530f;
        while (a < -3.14159265f) a += 6.28318530f;
        return a;
    }
};
\end{verbatim}

\textbf{用法（Usage）：}
\begin{verbatim}
    EKF_IMU ekf;
    ekf.dt = 0.002f;  // IMU 采样频率 500 Hz
    ekf.init();

    // 针对你的 IMU 进行调参（MPU6050 示例）：
    ekf.Q(0,0) = 0.001f;  ekf.Q(1,1) = 0.001f;  ekf.Q(2,2) = 0.001f;
    ekf.Q(3,3) = 1e-5f;   ekf.Q(4,4) = 1e-5f;   ekf.Q(5,5) = 1e-5f;
    ekf.R(0,0) = 0.1f;    ekf.R(1,1) = 0.1f;

    // 在 IMU 中断中（500 Hz）：
    ekf.update(gyro_x, gyro_y, gyro_z, accel_x, accel_y, accel_z);

    float roll_deg  = ekf.roll()  * 180.0f / M_PI;
    float pitch_deg = ekf.pitch() * 180.0f / M_PI;
    float yaw_deg   = ekf.yaw()   * 180.0f / M_PI;
\end{verbatim}

\subsection{TinyMPC 求解器（TinyMPC Solver）}

基于 Riccati 方法的最小化嵌入式 MPC 求解器（参见第 6 节理论部分）。可在微控制器上实时求解带盒式约束的 LQR 问题，基于 TinyMPC（Nguyen 等，2024）。

\begin{verbatim}
// tiny_mpc.hpp — 适用于嵌入式系统的基于 Riccati 的 MPC
// 基于：TinyMPC（Nguyen et al., 2024）
// 离线预计算 Riccati 递推，在线用 ADMM 求解。
#pragma once
#include "mat.hpp"

template<int NX, int NU, int N_HORIZON>
class TinyMPC {
public:
    using StateVec = Mat<NX, 1>;
    using InputVec = Mat<NU, 1>;
    using StateMat = Mat<NX, NX>;
    using InputMat = Mat<NX, NU>;
    using GainMat  = Mat<NU, NX>;

    // 系统模型（离散）
    StateMat Ad;
    InputMat Bd;

    // 代价矩阵
    StateMat Q;          // 状态代价
    Mat<NU,NU> R;        // 输入代价
    StateMat Q_terminal; // 终端代价（设为 DARE 解）

    // 约束
    InputVec u_min, u_max;
    StateVec x_min, x_max;

    // ADMM 参数
    float rho = 1.0f;        // 惩罚参数
    int   max_iter = 10;      // ADMM 迭代次数（典型值 5-20）

    // 预计算的 Riccati 增益（离线调用 precompute() 一次）
    GainMat  K_gains[N_HORIZON];    // 反馈增益
    StateMat P_mats[N_HORIZON + 1]; // 代价函数矩阵

    /// 预计算 Riccati 递推（调用一次，或在模型变化时重新调用）
    void precompute() {
        P_mats[N_HORIZON] = Q_terminal;
        for (int k = N_HORIZON - 1; k >= 0; --k) {
            // P = Q + A'*P_next*A - A'*P_next*B*(R+B'*P_next*B)^-1
            //     * B'*P_next*A
            auto P_next = P_mats[k + 1];
            auto BtP = Bd.T() * P_next;
            auto S = R + BtP * Bd;  // NU x NU
            auto S_inv = S.inverse();
            K_gains[k] = S_inv * BtP * Ad;  // NU x NX
            auto AtP = Ad.T() * P_next;
            P_mats[k] = Q + AtP * Ad - AtP * Bd * K_gains[k];
        }
    }

    /// 针对当前状态求解 MPC，返回最优首步输入
    InputVec solve(const StateVec& x0,
                   const StateVec x_ref[N_HORIZON]) {
        // 用无约束 LQR 展开初始化
        StateVec x_traj[N_HORIZON + 1];
        InputVec u_traj[N_HORIZON];
        x_traj[0] = x0;

        for (int k = 0; k < N_HORIZON; ++k) {
            StateVec dx = x_traj[k] - x_ref[k];
            u_traj[k] = (K_gains[k] * dx) * (-1.0f);
            // 截断（投影步骤）
            for (int i = 0; i < NU; ++i) {
                float v = u_traj[k](i, 0);
                v = v < u_min(i,0) ? u_min(i,0) :
                    (v > u_max(i,0) ? u_max(i,0) : v);
                u_traj[k](i, 0) = v;
            }
            x_traj[k + 1] = Ad * x_traj[k] + Bd * u_traj[k];
        }

        // ADMM 精化以满足约束
        InputVec z[N_HORIZON];      // 松弛变量
        InputVec lambda[N_HORIZON]; // 对偶变量
        for (int k = 0; k < N_HORIZON; ++k) {
            z[k] = u_traj[k];
            lambda[k] = InputVec::zeros();
        }

        for (int iter = 0; iter < max_iter; ++iter) {
            // 前向传播：带惩罚的无约束求解
            x_traj[0] = x0;
            for (int k = 0; k < N_HORIZON; ++k) {
                StateVec dx = x_traj[k] - x_ref[k];
                // 含 ADMM 惩罚的改进 Riccati
                InputVec u_unc = (K_gains[k] * dx) * (-1.0f);
                // ADMM 修正
                u_traj[k] = u_unc + (z[k] - lambda[k]) * (rho * 0.5f);
                x_traj[k+1] = Ad * x_traj[k] + Bd * u_traj[k];
            }
            // z 更新：投影到约束集
            for (int k = 0; k < N_HORIZON; ++k) {
                for (int i = 0; i < NU; ++i) {
                    float v = u_traj[k](i,0) + lambda[k](i,0);
                    v = v < u_min(i,0) ? u_min(i,0) :
                        (v > u_max(i,0) ? u_max(i,0) : v);
                    z[k](i, 0) = v;
                }
            }
            // 对偶变量更新
            for (int k = 0; k < N_HORIZON; ++k) {
                lambda[k] = lambda[k] + u_traj[k] - z[k];
            }
        }
        return u_traj[0];
    }
};
\end{verbatim}

\textbf{用法——底盘轨迹跟踪：}
\begin{verbatim}
    // 二维位置控制：状态=[x,y,vx,vy]，输入=[ax,ay]
    TinyMPC<4, 2, 10> mpc;  // 4 状态，2 输入，预测步长=10
    float dt = 0.01f;

    // 离散双积分器模型
    mpc.Ad = Mat<4,4>::identity();
    mpc.Ad(0,2) = dt; mpc.Ad(1,3) = dt;
    mpc.Bd(2,0) = dt; mpc.Bd(3,1) = dt;

    // 代价权重
    mpc.Q = Mat<4,4>::identity() * 10.0f;
    mpc.R = Mat<2,2>::identity() * 1.0f;
    mpc.Q_terminal = Mat<4,4>::identity() * 100.0f;

    // 约束：最大加速度 5 m/s^2
    mpc.u_min(0,0) = -5; mpc.u_min(1,0) = -5;
    mpc.u_max(0,0) =  5; mpc.u_max(1,0) =  5;

    mpc.precompute();  // 启动时调用一次

    // 在控制循环中：
    Mat<4,1> x_ref[10];  // 由路径规划器填充
    auto u = mpc.solve(current_state, x_ref);
    set_accel(u(0,0), u(1,0));
\end{verbatim}

\subsection{模块依赖关系图（Module Dependency Map）}

\begin{center}
\begin{tabular}{lll}
\toprule
\textbf{模块} & \textbf{文件} & \textbf{依赖} \\
\midrule
低通滤波器 & \texttt{lpf.hpp} & （无） \\
PID 控制器 & \texttt{pid.hpp} & （无） \\
高级 PID & \texttt{pid\_advanced.hpp} & （无） \\
给定值斜坡 & \texttt{ramp.hpp} & \texttt{lpf.hpp} \\
互补滤波器 & \texttt{complementary\_filter.hpp} & （无） \\
矩阵库 & \texttt{mat.hpp} & （无） \\
卡尔曼滤波器 & \texttt{kalman.hpp} & \texttt{mat.hpp} \\
EKF IMU（6自由度） & \texttt{ekf\_imu.hpp} & \texttt{mat.hpp} \\
LQR 控制器 & \texttt{lqr.hpp} & \texttt{mat.hpp} \\
TinyMPC 求解器 & \texttt{tiny\_mpc.hpp} & \texttt{mat.hpp} \\
串级 PID & \texttt{cascaded\_pid.hpp} & \texttt{pid.hpp} \\
\bottomrule
\end{tabular}
\end{center}

所有模块均采用 MIT 许可证。将其复制到项目的 \texttt{include/control/} 目录下，\texttt{\#include} 后即可使用。

\newpage

\vspace{1em}
\hrule
\vspace{1em}

\section*{结语}

如果你已读到这里，恭喜——你已经为设计竞赛机器人的控制系统打下了坚实的理论基础。以下是一张速查表，供你在不同场景下快速选用合适的工具：

\begin{center}
\begin{tabular}{ll}
\toprule
\textbf{场景} & \textbf{工具} \\
\midrule
传感器数据噪声大 & 低通滤波器（第 2 节） \\
电机实际特性是什么？ & 电机建模（第 3 节） \\
单电机转速/位置控制 & PID（第 4 节） \\
云台跟踪 & PID + 前馈（第 4 节） \\
平滑点对点运动 & 轨迹生成（第 4 节） \\
将 PID 从纸面移植到 MCU & Tustin 离散化（第 5 节） \\
平衡小车 / MIMO 系统 & LQR（第 6 节） \\
带约束的轨迹跟踪 & MPC（第 6 节） \\
噪声传感器的状态估计 & 卡尔曼滤波器（第 6 节） \\
IMU + 编码器融合 & 扩展卡尔曼滤波器（第 6 节） \\
机器人翻转/翻滚超过 $\pm 90$\textdegree{} & 四元数 EKF（第 6 节） \\
\bottomrule
\end{tabular}
\end{center}

请记住：理论是工具，而非目的。\textbf{最好的控制器是能在赛场上正常运行的那一个}，而不是在纸上看起来最优雅的那一个。从简单的方法（PID）开始，理解你的系统，只在必要时增加复杂度，并且始终、始终在真实硬件上进行测试。

祝你在下次比赛中好运。去造出令人惊叹的东西吧。

\end{document}
